\documentclass[11pt]{article}

\usepackage[margin=1in]{geometry}
\usepackage{amsmath,amssymb,amsthm,mathtools}
\usepackage{microtype}
\usepackage{enumitem}
\usepackage[hidelinks]{hyperref}
\usepackage{booktabs}
\usepackage{xcolor}

\newtheorem{theorem}{Theorem}[section]
\newtheorem{proposition}[theorem]{Proposition}
\newtheorem{corollary}[theorem]{Corollary}
\newtheorem{lemma}[theorem]{Lemma}
\theoremstyle{remark}
\newtheorem{remark}[theorem]{Remark}
\newtheorem{example}[theorem]{Example}

\newcommand{\R}{\mathbb{R}}
\newcommand{\N}{\mathbb{N}}
\newcommand{\Om}{\Omega}
\newcommand{\eps}{\varepsilon}
\newcommand{\Lp}{L^p}
\newcommand{\Wp}{W^{1,p}}
\newcommand{\ellinf}{\ell_\infty}
\newcommand{\cnull}{c_0}
\newcommand{\rhoMin}[1]{\rho_{#1}}
\newcommand{\restr}[2]{#1\vert_{#2}}
\newcommand{\Wgraph}{\mathcal{W}^{1,p}}
\newcommand{\dist}{\operatorname{dist}}

\title{Banach-Valued Sobolev Extension Sets:\\
Verification, Current Status, and Further Partial Results}
\author{Research note prepared from the supplied partial solution}
\date{June 19, 2026}

\begin{document}
\maketitle

\begin{abstract}
The supplied finite-dimensional argument is correct.  This note gives a
fully quantitative version: if a scalar Newton--Sobolev extension map has
bound $C_E$, then an $n$-dimensional target $V$ has an extension map with
bound at most
\[
  2^{1-1/p} n^{1/p} C_E\,\|A\|\,\|A^{-1}\|,
  \qquad A:V\longrightarrow \ell_\infty^n.
\]
The proof requires no Poincare inequality and works for $1\le p<\infty$.
A literature search did not locate a complete positive solution or a
counterexample to Question 1.1 of Garcia-Bravo--Ikonen--Zhu.  In fact, a
May 2026 paper by the same authors explicitly states that even the Euclidean
case $1\le p<n$ remains open.

Two further infinite-dimensional positive results are established here in
the Euclidean setting.  If a scalar extension domain has a bounded linear
scalar extension operator, then it is an extension domain for targets
isomorphic to complemented subspaces of an $L^p$-space.  In particular this
covers $L^p$-targets, $\ell^p$, and arbitrary Hilbert spaces.  Finally, an
explicit $c_0$-valued Lipschitz map on an interval shows why finite-rank
approximation does not by itself pass the finite-dimensional argument to
$c_0$: canonical truncations converge uniformly but not in Newton--Sobolev
norm, and no finite-dimensional-range maps can converge to this example in
that norm.
\end{abstract}

\begin{center}
\fbox{\parbox{0.91\linewidth}{\textbf{Bottom line.}
The partial theorem is valid after minor technical cleanup.  No full answer
or counterexample was found.  The general problem remains open as of the
search date, including the Euclidean range $1<p<n$.}}
\end{center}

\section{The question and its current status}

Let $1<p<\infty$, let $Z$ be a $p$-PI metric measure space, and let
$\Om\subset Z$.  Question 1.1 of Garcia-Bravo--Ikonen--Zhu asks whether every
real-valued $\Wp$-extension set is automatically a $V$-valued
$\Wp$-extension set for every Banach space $V$ \cite{GBIZ23}.
Their definition does not require an extension map to be linear.

The 2023 paper proves several substantial positive cases.  Among them:
\begin{itemize}[leftmargin=2em,itemsep=2pt]
  \item a $c_0$-valued extension set is an extension set for every Banach
  target;
  \item measure density together with a weak $(1,p)$-Poincare inequality up
  to some scale gives extension for every Banach target;
  \item Euclidean $\Wp$-extension domains have the desired target-independent
  extension property when $p\ge n$;
  \item planar Jordan extension domains have it for all $1<p<\infty$.
\end{itemize}

The most informative later source found in the search is the May 29, 2026
revision of \emph{On infinity thick quasiconvexity and applications}, by the
same three authors.  It says explicitly that for domains in $\R^n$ the
equivalence has been verified for $p\in[n,\infty)$, while
$p\in[1,n)$ remains open \cite{GBIZ26}.  Consequently, a complete answer to
the original general $p$-PI question would in particular settle an open
Euclidean subcase.  No later full theorem or counterexample was located in
searches by title, authors, citations, and the phrases ``Banach-valued
Sobolev extension set/domain'' through June 19, 2026.

\section{Conventions}

We use the Newton--Sobolev convention of \cite{GBIZ23}.  For a Banach-valued
map $u$ on a metric measure space $X$, write $\rhoMin{u}$ for its minimal
$p$-weak upper gradient and set
\[
  \|u\|_{\Wp(X;V)}
  :=\|u\|_{\Lp(X;V)}+\|\rhoMin{u}\|_{\Lp(X)}.
\]
If $T:V\to W$ is bounded and linear, then
\[
  \|T\circ u\|_{\Lp(X;W)}\le \|T\|\,\|u\|_{\Lp(X;V)},
  \qquad
  \rhoMin{T\circ u}\le \|T\|\rhoMin{u}
  \quad\text{a.e.}
\]
The same upper-gradient conclusion holds for Lipschitz postcomposition.

An ``extension map of bound $C_E$'' means a map $E$ satisfying
$\restr{Eu}{\Om}=u$ and $\|Eu\|\le C_E\|u\|$.  If $E$ is nonlinear, the
phrase ``operator norm'' should be understood only as this Lipschitz-type
bound at the origin.

\section{Verification of the finite-dimensional result}

\begin{theorem}[Quantitative finite-dimensional extension]
\label{thm:finite}
Let $1\le p<\infty$, let $Z$ be a metric measure space, and let
$\Om\subset Z$ be measurable.  Suppose there is a scalar extension map
\[
  E_{\R}:\Wp(\Om)\longrightarrow \Wp(Z)
\]
with bound $C_E$.  Let $V$ be $n$-dimensional and let
$A:V\to\ell_\infty^n$ be a linear isomorphism.  Then there is a $V$-valued
extension map $E_V$ satisfying
\begin{equation}
\label{eq:finiteconstant}
  \|E_Vu\|_{\Wp(Z;V)}
  \le 2^{1-1/p} n^{1/p} C_E
       \|A\|\,\|A^{-1}\|\,
       \|u\|_{\Wp(\Om;V)}.
\end{equation}
If $E_{\R}$ is linear, then $E_V$ is linear.
\end{theorem}

\begin{proof}
We first take $V=\ell_\infty^n$.  Write
$u=(u_1,\dots,u_n)$ and put
\[
  h_i=E_{\R}u_i,
  \qquad
  H=(h_1,\dots,h_n).
\]
Each coordinate projection is $1$-Lipschitz, hence
$u_i\in\Wp(\Om)$ and $\rhoMin{u_i}\le\rhoMin{u}$ almost everywhere.  In
particular,
\begin{equation}
\label{eq:coordinput}
  \|u_i\|_{\Wp(\Om)}
  \le \|u\|_{\Wp(\Om;\ell_\infty^n)}.
\end{equation}

Let $g_i=\rhoMin{h_i}$ and set $g=\max_{1\le i\le n}g_i$.  Outside the finite
union of the exceptional $p$-negligible curve families for the $g_i$, every
rectifiable curve $\gamma$ satisfies
\[
 \begin{split}
  \|H(\gamma(a))-H(\gamma(b))\|_{\ell_\infty^n}
  &=\max_i |h_i(\gamma(a))-h_i(\gamma(b))|\\
  &\le \int_\gamma g\,ds.
 \end{split}
\]
Thus $g$ is a $p$-weak upper gradient of $H$.  Strong measurability is
automatic because the target is finite-dimensional, and the coordinate
restriction identities imply $\restr{H}{\Om}=u$ almost everywhere.

Put
\[
  a_i=\|h_i\|_{\Lp(Z)},
  \qquad b_i=\|g_i\|_{\Lp(Z)}.
\]
Then
\[
 \begin{split}
  \|H\|_{\Wp(Z;\ell_\infty^n)}
  &\le \Big(\sum_i a_i^p\Big)^{1/p}
       +\Big(\sum_i b_i^p\Big)^{1/p}\\
  &\le 2^{1-1/p}
       \Big(\sum_i(a_i^p+b_i^p)\Big)^{1/p}\\
  &\le 2^{1-1/p}
       \Big(\sum_i(a_i+b_i)^p\Big)^{1/p}\\
  &\le 2^{1-1/p}C_E
       \Big(\sum_i\|u_i\|_{\Wp(\Om)}^p\Big)^{1/p}\\
  &\le 2^{1-1/p}n^{1/p}C_E
       \|u\|_{\Wp(\Om;\ell_\infty^n)},
 \end{split}
\]
where the last inequality uses \eqref{eq:coordinput}.

For general $V$, define
\[
  E_Vu=A^{-1}E_{\ell_\infty^n}(Au).
\]
The bounded-linear postcomposition property gives
\[
  \|Au\|_{\Wp(\Om;\ell_\infty^n)}
  \le \|A\|\,\|u\|_{\Wp(\Om;V)}
\]
and the analogous estimate for $A^{-1}$ on $Z$.  This proves
\eqref{eq:finiteconstant}.  Linearity is inherited coordinatewise when the
scalar extension map is linear.
\end{proof}

\begin{corollary}[A dimension-only bound]
Every $n$-dimensional Banach space admits a linear isomorphism
$A:V\to\ell_\infty^n$ with
$\|A\|\,\|A^{-1}\|\le n$.  Hence Theorem \ref{thm:finite} always gives the
bound
\[
  2^{1-1/p}n^{1+1/p}C_E.
\]
\end{corollary}

\begin{proof}
Choose an Auerbach basis $(v_i,v_i^*)_{i=1}^n$, so that
$\|v_i\|=\|v_i^*\|=1$ and $v_i^*(v_j)=\delta_{ij}$.  The coordinate map
$A(v)=(v_i^*(v))_{i=1}^n$ has norm at most one, while
$\|v\|\le\sum_i|v_i^*(v)|\le n\|Av\|_\infty$.
\end{proof}

\subsection*{What was right, and what needed tightening}

The supplied packet had the correct central idea.  The following are the
only substantive cleanups.
\begin{enumerate}[leftmargin=2em,itemsep=3pt]
  \item No PI, doubling, completeness, or measure-density hypothesis is
  needed for the finite-dimensional statement.
  \item The coordinate estimate is exact for the stated Newton--Sobolev
  norm; it is not merely true ``up to an equivalent norm.''
  \item One should explicitly take the finite union of the exceptional
  curve families.
  \item The extension is linear exactly when the original scalar extension
  is linear.
  \item Formula \eqref{eq:finiteconstant} records the quantitative loss.
\end{enumerate}
The same finite-coordinate mechanism also appears in Lemma 4.5 of
\cite{GBIZ23}, there in a density argument rather than as an extension
theorem.

\section{A further infinite-dimensional positive result in Euclidean space}

The next result does not solve the open problem, but it goes beyond finite
dimension.  Its key point is that the Sobolev graph norm tensorizes exactly
with an $L^p$ target when the target exponent matches the Sobolev exponent.

For an open $D\subset\R^n$ and a Banach space $V$ with the Radon--Nikodym
property, use the equivalent classical graph norm
\begin{equation}
\label{eq:graphnorm}
  \|u\|_{\Wgraph(D;V)}
  :=\left(
       \|u\|_{\Lp(D;V)}^p+
       \sum_{j=1}^n\|\partial_j u\|_{\Lp(D;V)}^p
     \right)^{1/p}.
\end{equation}
On Euclidean open sets, the Newton--Sobolev and Reshetnyak spaces agree for
all Banach targets, and the classical Bochner--Sobolev space agrees with
them precisely for targets having the Radon--Nikodym property
\cite{CaamanoEtAl}.

\begin{lemma}[Sobolev--Fubini identification]
\label{lem:fubini}
Let $(S,\nu)$ be a $\sigma$-finite measure space and $D\subset\R^n$ be open.
Under the identification $U(s)(x)=u(x)(s)$,
\[
  \Wgraph(D;L^p(S,\nu))
  \cong L^p\bigl(S,\nu;\Wgraph(D)\bigr)
\]
isometrically for the norm \eqref{eq:graphnorm}.
\end{lemma}

\begin{proof}
For a Bochner--Sobolev map $u$, Fubini's theorem and the distributional
definition of the weak partial derivatives show that for almost every
$s$, the slice $U(s)$ lies in $\Wgraph(D)$ and
$\partial_j U(s)=(\partial_j u)(\cdot)(s)$.  Conversely, the same
calculation against tensor-product test functions gives the weak derivatives
of $u$ from those of the slices.  Finally,
\[
 \begin{split}
  \|u\|_{\Wgraph(D;L^p(S))}^p
  &=\int_S\left(
       \|U(s)\|_{L^p(D)}^p+
       \sum_{j=1}^n\|\partial_jU(s)\|_{L^p(D)}^p
     \right)d\nu(s),
 \end{split}
\]
which is the asserted isometry.
\end{proof}

\begin{theorem}[$L^p$-target tensorization]
\label{thm:Lptarget}
Let $1<p<\infty$, let $\Om\subset\R^n$ be open, and suppose there is a
bounded linear scalar extension operator
\[
  E:\Wgraph(\Om)\longrightarrow \Wgraph(\R^n).
\]
Then, for every $\sigma$-finite $(S,\nu)$, there is a bounded linear extension
operator
\[
  E_{L^p(S)}:\Wgraph(\Om;L^p(S,\nu))
  \longrightarrow \Wgraph(\R^n;L^p(S,\nu))
\]
with the same operator bound when both spaces use the graph norm
\eqref{eq:graphnorm}.
\end{theorem}

\begin{proof}
By Lemma \ref{lem:fubini}, an $L^p(S)$-valued Sobolev map corresponds to a
Bochner map
\[
  U\in L^p\bigl(S;\Wgraph(\Om)\bigr).
\]
Every bounded linear map between Banach spaces acts pointwise, with the same
norm, on the corresponding Bochner $L^p$ spaces.  Define
\[
  \widetilde E U(s):=E(U(s)).
\]
Then
\[
  \|\widetilde E U\|_{L^p(S;\Wgraph(\R^n))}
  \le \|E\|\,
      \|U\|_{L^p(S;\Wgraph(\Om))}.
\]
Transporting $\widetilde E U$ back through Lemma \ref{lem:fubini} gives the
desired extension.  The scalar restriction identity holds for almost every
$s$, hence the resulting $L^p(S)$-valued map restricts to the original map.
\end{proof}

\begin{corollary}[Complemented subspaces of $L^p$]
\label{cor:complemented}
Under the hypotheses of Theorem \ref{thm:Lptarget}, let $V$ be isomorphic to
a complemented subspace of some $L^p(S,\nu)$.  Then $\Om$ is a
$V$-valued extension domain.  Quantitatively, if
$J:V\to L^p(S,\nu)$ is an isomorphic embedding and
$Q:L^p(S,\nu)\to J(V)$ is a bounded projection, the extension bound is at
most
\[
  \|J^{-1}\|\,\|Q\|\,\|E\|\,\|J\|.
\]
\end{corollary}

\begin{proof}
Apply Theorem \ref{thm:Lptarget} to $J\circ u$, then postcompose the extension
with $J^{-1}Q$.
\end{proof}

\begin{corollary}[Concrete targets]
\label{cor:concrete}
Let $1<p<\infty$ and let $\Om\subset\R^n$ be a scalar $\Wp$-extension domain.
Then $\Om$ is a Newton--Sobolev extension domain for each of the following
infinite-dimensional targets:
\[
  L^p(S,\nu),\qquad \ell^p,\qquad\text{and every Hilbert space }H.
\]
More generally, it works for every Banach space isomorphic to a complemented
subspace of an $L^p$-space.
\end{corollary}

\begin{proof}
A Euclidean scalar extension domain satisfies the measure-density condition
\cite{HKT08b}.  The linearization theorem in \cite{GBIZ23} therefore provides
a bounded linear scalar extension operator, so Theorem
\ref{thm:Lptarget} and Corollary \ref{cor:complemented} apply.

It remains only to recall why Hilbert spaces are covered.  Let
$H=\ell^2(I)$ and let $(r_i)_{i\in I}$ be independent Rademacher functions on
$\{-1,1\}^I$ with product probability measure.  Khintchine's inequality
gives an isomorphic embedding
\[
  J:\ell^2(I)\longrightarrow L^p(\{-1,1\}^I),
  \qquad Ja=\sum_{i\in I}a_i r_i.
\]
For $f\in L^p$, set $c_i=\int f r_i$.  By duality and Khintchine's inequality
for $p'=p/(p-1)$,
\[
  \|(c_i)\|_{\ell^2(I)}\le C_{p'}\|f\|_{L^p}.
\]
Thus
\[
  Qf:=\sum_{i\in I}c_i r_i
\]
defines a bounded projection onto the Rademacher span.  Hence $H$ is
isomorphic to a complemented subspace of an $L^p$-space.  Such targets are
reflexive, hence have the Radon--Nikodym property, so the classical and
Newton--Sobolev formulations agree by \cite{CaamanoEtAl}.
\end{proof}

\begin{remark}
Corollary \ref{cor:concrete} is not a solution for arbitrary Banach targets.
The tensor extension of a general bounded scalar operator need not be bounded
on $L^p(\cdot;V)$ for an arbitrary $V$; complemented $L^p$ structure is the
precise mechanism used above.
\end{remark}

\section{Why finite-rank approximation does not finish the problem}

The source paper proves energy-density for targets with the metric
approximation property but deliberately does not claim norm-density.  The
following elementary example shows a sharp obstruction.

\begin{proposition}[$c_0$ truncations can fail completely]
\label{prop:c0}
Let $I=(0,1)$ and define
\[
  u(t)=\left(\frac{\sin(kt)}{k}\right)_{k=1}^{\infty}\in c_0.
\]
For every $1\le p<\infty$:
\begin{enumerate}[label=\textup{(\roman*)},leftmargin=2em]
  \item $u$ is $1$-Lipschitz, hence belongs to the Newton--Sobolev space
  $\Wp(I;c_0)$;
  \item if $P_N$ is the first-$N$ coordinate projection, then
  $P_Nu\to u$ uniformly (and hence in $L^p$), but
  \[
    \rhoMin{u-P_Nu}=1\quad\text{a.e. on }I;
  \]
  in particular $P_Nu$ does not converge to $u$ in $\Wp(I;c_0)$;
  \item $u$ is not in the classical Bochner--Sobolev space
  $\Wp_{\mathrm{Bochner}}(I;c_0)$;
  \item no sequence of Newton--Sobolev maps whose essential ranges lie in
  finite-dimensional subspaces of $c_0$ can converge to $u$ in Newton--Sobolev
  norm.
\end{enumerate}
\end{proposition}

\begin{proof}
Every coordinate $t\mapsto \sin(kt)/k$ is $1$-Lipschitz, so the supremum
norm gives
$\|u(t)-u(s)\|_{c_0}\le |t-s|$.  Also
$|\sin(kt)|/k\to0$, so the image lies in $c_0$.

For $v_N=u-P_Nu$,
\[
  \|v_N\|_{L^\infty(I;c_0)}\le\frac1{N+1}.
\]
On the other hand, for every $t$,
\[
  \sup_{k>N}|\cos(kt)|=1.
\]
Indeed, this is immediate when $t/(2\pi)$ is rational by taking multiples of
a denominator, and follows from density of the irrational rotation when
$t/(2\pi)$ is irrational.  The metric derivative therefore satisfies
\[
 \begin{split}
  \operatorname{md}v_N(t)
  &=\lim_{h\to0}
    \frac{\|v_N(t+h)-v_N(t)\|_{c_0}}{|h|}=1.
 \end{split}
\]
The upper bound follows from $1$-Lipschitz continuity, while the lower bound
is obtained by fixing a coordinate with $|\cos(kt)|$ arbitrarily close to
one.  On an interval, the minimal weak upper gradient of a Lipschitz
metric-valued map equals its metric derivative almost everywhere, proving
(ii).

If a Bochner derivative $u'(t)\in c_0$ existed, every coordinate projection
would give
\[
  u'(t)=(\cos(kt))_{k=1}^{\infty}.
\]
This sequence is not in $c_0$ for any $t$, so $u$ is nowhere Bochner
differentiable.  This proves (iii); it is the standard example recorded in
\cite{CaamanoEtAl}.

Finally, on Euclidean open sets the Newton--Sobolev and Reshetnyak spaces
coincide, while the classical Bochner--Sobolev space is a closed subspace of
the Reshetnyak space \cite{CaamanoEtAl}.  Every finite-dimensional-range
Newton--Sobolev map is Bochner--Sobolev, because finite-dimensional targets
have the Radon--Nikodym property.  Norm convergence of such maps to $u$ would
therefore force $u$ into that closed classical subspace, contradicting (iii).
\end{proof}

\begin{remark}
The proposition does not produce a counterexample to the extension question:
$I$ itself is already a PI space and $u$ needs no extension.  It rules out a
natural proof strategy.  One cannot pass from the finite-dimensional theorem
to $c_0$ merely by truncating the target and taking a Newton--Sobolev norm
limit.
\end{remark}

\section{Strategies tested against the remaining problem}

The following approaches were examined; each isolates a genuine obstruction.

\paragraph{Coordinatewise extension to $c_0$ or $\ell^\infty$.}
Finite coordinates give Theorem \ref{thm:finite}, but the bound grows with
the number of coordinates.  For infinitely many coordinates, scalar bounds
do not ensure that the assembled map is Bochner measurable, $L^p$-bounded,
or controlled by one common upper gradient.  Proposition \ref{prop:c0}
shows that finite-coordinate limits need not converge in Sobolev norm.

\paragraph{Metric approximation property.}
Finite-rank operators approximate compact subsets of the target and yield
the energy-density theorem in \cite{GBIZ23}.  Energy convergence is weaker
than Newton--Sobolev norm convergence; Proposition \ref{prop:c0} shows that
the gap is not a technicality.

\paragraph{Scalarization by the dual unit ball.}
For every $v^*\in B_{V^*}$, the scalar function $v^*\circ u$ extends.  The
problem is to choose these scalar extensions compatibly.  Uniform scalar
norm bounds do not produce a single pointwise $L^p$ upper gradient valid for
all $v^*$, nor do they guarantee that the values come from one element of
$V$ rather than merely from $V^{**}$ or a product space.

\paragraph{Linear tensor extension.}
This succeeds for complemented subspaces of $L^p$ by Section 4, and for
Hilbert targets through the Rademacher complement.  For arbitrary $V$, a
bounded scalar operator on an $L^p$-type graph subspace need not have a
bounded amplification by $\operatorname{Id}_V$.  Thus ordinary boundedness of
the scalar extension operator is not enough for abstract tensorization.

\paragraph{A nonlinear-Poincare or expander counterexample.}
One might try to encode a family with good scalar Poincare inequalities but
bad $c_0$-valued inequalities.  Classical expanders provide that functional-
analytic contrast, but an unbounded expander family is not uniformly
doubling.  Since every subset of a doubling metric space is metrically
doubling, this mechanism cannot be inserted at uniform scale into a fixed
PI ambient space without losing the hypotheses.  No valid PI realization
was found.

\paragraph{Removable sets and zero-capacity bottlenecks.}
If the removed set has zero $p$-capacity, the family of curves meeting it is
$p$-negligible, a target-independent fact.  Such examples therefore tend to
extend all metric or Banach targets, rather than separate scalar and vector
extension.  Positive-capacity bottlenecks usually already obstruct scalar
extension.

These failures are consistent with the formulation in \cite{GBIZ23}: in the
measure-dense, path-almost-open setting, the open issue can be reduced to a
closed-image/completeness question for the $c_0$-valued local
Hajlasz--Sobolev space.  None of the arguments above establishes or refutes
that completeness.

\section{Conclusion}

The supplied partial result is mathematically sound and can be strengthened
to the quantitative Theorem \ref{thm:finite}.  The literature does not
currently supply a full answer or counterexample; the same authors still
identify the Euclidean range $1<p<n$ as open in May 2026.  The tensorization
argument above adds a genuine infinite-dimensional class in Euclidean
space---all targets complemented in an $L^p$-space, including $L^p$,
$\ell^p$, and Hilbert spaces---but arbitrary Banach targets remain beyond
it.  The $c_0$ example explains why the most direct finite-rank limiting
argument cannot bridge the remaining gap.

\begin{thebibliography}{99}

\bibitem{GBIZ23}
M. Garcia-Bravo, T. Ikonen, and Z. Zhu,
\emph{Extensions and approximations of Banach-valued Sobolev functions},
Math. Z. \textbf{305} (2023), no. 4, Paper No. 67, 50 pp.;
arXiv:2208.12594.

\bibitem{GBIZ26}
M. Garcia-Bravo, T. Ikonen, and Z. Zhu,
\emph{On infinity thick quasiconvexity and applications},
arXiv:2509.01194v2, May 29, 2026.

\bibitem{CaamanoEtAl}
I. Caamano, J. A. Jaramillo, A. Prieto, and A. Ruiz de Alarcon,
\emph{Sobolev spaces of vector-valued functions},
arXiv:2008.03040; see in particular Theorems 4.3 and 4.6.

\bibitem{HKT08b}
P. Hajlasz, P. Koskela, and H. Tuominen,
\emph{Sobolev embeddings, extensions and measure density condition},
J. Funct. Anal. \textbf{254} (2008), 1217--1234.

\bibitem{HKST15}
J. Heinonen, P. Koskela, N. Shanmugalingam, and J. T. Tyson,
\emph{Sobolev Spaces on Metric Measure Spaces: An Approach Based on Upper
Gradients}, Cambridge University Press, 2015.

\end{thebibliography}

\end{document}
